\pagestyle{empty}
\tight
\begin{tblr}{width=\textwidth,colspec={|Q[c,m,1.9cm]|X[l,m]|},hlines,vlines,hspan=minimal,row{1}={bg=headblue},rowsep=1pt,colsep=3pt}
\SetCell{c,m}\hfil\logo\hfil & \textbf{POLITEKNIK NEGERI MALANG}\par\textbf{JURUSAN TEKNOLOGI INFORMASI}\par\textbf{PROGRAM STUDI: D4 TEKNIK INFORMATIKA} \\
\SetCell[c=2]{c} \textbf{RENCANA PEMBELAJARAN SEMESTER (RPS)} & \\
\end{tblr}
\begin{longtblr}{width=\textwidth,colspec={|Q[l,m,1.9cm]|X[0.85,l,m]|X[1.45,l,m]|X[0.75,l,m]|X[1.15,l,m]|},hlines,vlines,hspan=minimal,rowsep=1pt,colsep=2pt,row{1,3}={bg=headgray,font=\bfseries}}
MATA KULIAH & KODE & BOBOT (SKS)/jam & SEMESTER & TGL. PENYUSUNAN \\
Pemrograman Web Lanjut & RTI254007 & 3 SKS/6 jam & 4 & 06 Februari 2025 \\
OTORISASI & Dosen Pengembang RPS & Koordinator RMK & \SetCell[c=2]{l} Ka PRODI &  \\
 & Moch. Zawaruddin Abdullah, S.ST., M.Kom. &  & \SetCell[c=2]{l}{\vspace{8mm}\par Dr. Ely Setyo Astuti, S.T., M.T.} &  \\
\SetCell[r=12]{l} Capaian Pembelajaran (CP) & \SetCell[c=4]{l,bg=headgray,font=\bfseries} Capaian Pembelajaran Lulusan Program Studi (CPL-Prodi) &  &  &  \\
 & CPL02 & \SetCell[c=3]{l} Mampu menerapkan prinsip rekayasa sistem dan perangkat lunak untuk menghasilkan solusi aplikasi multi-platform sesuai kebutuhan. &  &  \\
 & CPL06 & \SetCell[c=3]{l} Mampu merancang, mengimplementasikan, menguji, melakukan deployment, memelihara, serta menjamin mutu perangkat lunak yang menjawab permasalahan dan memenuhi kebutuhan pengguna. &  &  \\
 & CPL10 & \SetCell[c=3]{l} Mampu menganalisis persoalan kompleks dengan wawasan multidisiplin untuk merumuskan solusi informatika atau ilmu komputer. &  &  \\
 & \SetCell[c=4]{l,bg=headgray,font=\bfseries} Capaian Pembelajaran mata kuliah (CPL-MK) &  &  &  \\
 & CPMK0209 & \SetCell[c=3]{l} Mahasiswa mampu memilih dan menerapkan teknologi pengembangan berbasis platform untuk membangun aplikasi sesuai kebutuhan. &  &  \\
 & CPMK0603 & \SetCell[c=3]{l} Mahasiswa mampu mengimplementasikan dan melakukan deployment aplikasi web, mobile, atau framework sesuai kebutuhan pengguna. &  &  \\
 & CPMK1006 & \SetCell[c=3]{l} Mahasiswa mampu mengevaluasi dan mengembangkan aplikasi web atau mobile berdasarkan kebutuhan, kualitas implementasi, dan pengalaman pengguna. &  &  \\
 & \SetCell[c=4]{l,bg=headgray,font=\bfseries} Kemampuan akhir tiap tahapan belajar (Sub-CPMK) &  &  &  \\
 & SCPMK0209-03401 & \SetCell[c=3]{l} Mahasiswa mampu menerapkan konsep framework web, MVC, routing, controller, view, dan version control. &  &  \\
 & SCPMK0603-03402 & \SetCell[c=3]{l} Mahasiswa mampu mengimplementasikan migration, ORM, Blade, layout, validasi, authentication, authorization, import/export, dan API pada aplikasi Laravel. &  &  \\
 & SCPMK1006-03403 & \SetCell[c=3]{l} Mahasiswa mampu mengevaluasi kualitas aplikasi Laravel berdasarkan fungsi, keamanan, API, dan kebutuhan pengguna. &  &  \\
\SetCell[r=2]{l} Penilaian & \SetCell[c=4]{l} *Nilai CPMK didapat dari agregasi nilai SUB-CPMK yang dibebankan &  &  &  \\
 & \SetCell[c=4]{l}{\tiny\begin{tblr}{width=\linewidth,colspec={|Q[c,m,0.1\linewidth]|Q[c,m,0.16\linewidth]|X[l,m]|X[l,m]|X[l,m]|X[l,m]|X[l,m]|Q[c,m,0.08\linewidth]|},hlines,vlines,hspan=minimal,rowsep=1pt,colsep=0.7pt,row{1,2}={bg=headgray,font=\bfseries}}
Kode CPMK & Kode SCPMK & \SetCell[c=5]{c} Bobot per Bentuk Penilaian & & & & & Total Bobot \\
 & & Tugas & UTS & Kuis & PBL & UAS & \\
CPMK0209 & SCPMK0209-03401 & 10\%: instalasi, routing, controller, view & 10\%: CRUD Laravel & & 3.33\%: proyek Laravel & & 23.33\% \\
CPMK0603 & SCPMK0603-03402 & 35\%: migration, ORM, Blade, auth, import/export, API & 10\%: CRUD Laravel & & 5.34\%: integrasi proyek & & 50.34\% \\
CPMK1006 & SCPMK1006-03403 & 8\%: testing, review kode, deployment & & 5\%: review kualitas Laravel & 5\%: proyek Laravel & 8.33\%: validasi proyek & 26.33\% \\
\SetCell[c=7]{r}\textbf{Total Per Penilaian} & & & & & & & \textbf{100\%} \\
\end{tblr}} &  &  &  \\
Diskripsi Singkat Mata Kuliah & \SetCell[c=4]{l} Mata kuliah Pemrograman Web Lanjut membahas pengembangan aplikasi web modern menggunakan Laravel, mencakup MVC, routing, migration, ORM, Blade, validasi, authentication, authorization, import/export, API, testing, dan evaluasi kualitas aplikasi. &  &  &  \\
Bahan Kajian / Materi Pembelajaran & \SetCell[c=4]{l}\begin{enumerate}\item Laravel, MVC, routing, controller, view, dan Git\item Migration, Eloquent ORM, Blade, layout, dan validasi\item Authentication, authorization, import/export, API, testing, dan proyek Laravel\end{enumerate} &  &  &  \\
\SetCell[r=2]{l} Pustaka & \SetCell[c=4]{l} \textbf{Utama:}\begin{enumerate}\item Azamuddin, M. and Mukhlasin, H. 2019. Laravel the PHP Framework for Web Artisans.\item Laravel Documentation. 2024. https://laravel.com/docs/10.x.\end{enumerate} &  &  &  \\
 & \SetCell[c=4]{l} \textbf{Pendukung:} Materi LMS, dokumentasi resmi perangkat lunak, dan jobsheet praktikum. &  &  &  \\
Media Pembelajaran & \SetCell[c=2]{l} \textbf{Software:} PHP, MySQL, Laravel, Composer, Git, Laragon, Visual Studio Code. &  & \SetCell[c=2]{l} \textbf{Hardware:} PC/laptop, proyektor. &  \\
Nama Dosen Pengampu & \SetCell[c=4]{l}\begin{enumerate}\item Farid Angga Pribadi, S.Kom., M.Kom.\item Moch. Zawaruddin Abdullah, S.ST., M.Kom.\item Dimas Wahyu Wibowo, S.T., M.T.\item Habibie Ed Dien, S.Kom., M.T.\item Endah Septa Sintiya, S.Pd., M.Kom.\item Muhammad Unggul Pamenang, S.St., M.T.\end{enumerate} &  &  &  \\
Matakuliah Syarat & \SetCell[c=4]{l} Desain dan Pemrograman Web &  &  &  \\
\end{longtblr}
\begin{longtblr}{width=\textwidth,colspec={|Q[c,m,0.06\textwidth]|X[1.35,l,m]|X[1.08,l,m]|X[1.16,l,m]|X[0.7,l,m]|X[1.24,l,m]|X[1.06,l,m]|X[0.98,l,m]|Q[c,m,0.05\textwidth]|},hlines,vlines,hspan=minimal,rowhead=2,row{1,2}={bg=headgreen,font=\bfseries},rowsep=1pt,colsep=1.5pt}
Mgg.\par Ke & Kemampuan Akhir Yang Direncanakan (Sub-CP-MK) & Bahan kajian (Materi Pembelajaran) & Bentuk dan Metode Pembelajaran & Estimasi Waktu & Pengalaman Belajar Mahasiswa & Kriteria \& Bentuk Penilaian & Indikator Penilaian & Bobot Penilaian (\%) \\
(1) & (2) & (3) & (4) & (5) & (6) & (7) & (8) & (9) \\
1 & SCPMK0209-03401 & Web framework dan instalasi Laravel. & Modalitas: Blended Learning. Bentuk: Praktikum/PBL. Metode: hands-on practice, mentoring, code review, demo, dan presentasi. & 1 x 2 x 50' tatap muka; 1 x 2 x 50' tugas/praktik mandiri. & Mahasiswa mengerjakan aktivitas terarah untuk web framework dan instalasi laravel dan menunjukkan evidence hasil belajar. & Bentuk penilaian: Web framework dan instalasi Laravel. Mengacu rubrik RTM. & Ketepatan konsep; kualitas artefak/praktik; validasi dan dokumentasi. & 5 \\
2 & SCPMK0209-03401 & Routing, controller, view. & Modalitas: Blended Learning. Bentuk: Praktikum/PBL. Metode: hands-on practice, mentoring, code review, demo, dan presentasi. & 1 x 2 x 50' tatap muka; 1 x 2 x 50' tugas/praktik mandiri. & Mahasiswa mengerjakan aktivitas terarah untuk routing, controller, view dan menunjukkan evidence hasil belajar. & Bentuk penilaian: Routing, controller, view. Mengacu rubrik RTM. & Ketepatan konsep; kualitas artefak/praktik; validasi dan dokumentasi. & 5 \\
3 & SCPMK0603-03402 & Migration. & Modalitas: Blended Learning. Bentuk: Praktikum/PBL. Metode: hands-on practice, mentoring, code review, demo, dan presentasi. & 1 x 2 x 50' tatap muka; 1 x 2 x 50' tugas/praktik mandiri. & Mahasiswa mengerjakan aktivitas terarah untuk migration dan menunjukkan evidence hasil belajar. & Bentuk penilaian: Migration. Mengacu rubrik RTM. & Ketepatan konsep; kualitas artefak/praktik; validasi dan dokumentasi. & 5 \\
4 & SCPMK0603-03402 & Eloquent ORM. & Modalitas: Blended Learning. Bentuk: Praktikum/PBL. Metode: hands-on practice, mentoring, code review, demo, dan presentasi. & 1 x 2 x 50' tatap muka; 1 x 2 x 50' tugas/praktik mandiri. & Mahasiswa mengerjakan aktivitas terarah untuk eloquent orm dan menunjukkan evidence hasil belajar. & Bentuk penilaian: Eloquent ORM. Mengacu rubrik RTM. & Ketepatan konsep; kualitas artefak/praktik; validasi dan dokumentasi. & 5 \\
5 & SCPMK0603-03402 & Blade, layout, datatable. & Modalitas: Blended Learning. Bentuk: Praktikum/PBL. Metode: hands-on practice, mentoring, code review, demo, dan presentasi. & 1 x 2 x 50' tatap muka; 1 x 2 x 50' tugas/praktik mandiri. & Mahasiswa mengerjakan aktivitas terarah untuk blade, layout, datatable dan menunjukkan evidence hasil belajar. & Bentuk penilaian: Blade, layout, datatable. Mengacu rubrik RTM. & Ketepatan konsep; kualitas artefak/praktik; validasi dan dokumentasi. & 5 \\
6 & SCPMK0603-03402 & Form dan validasi. & Modalitas: Blended Learning. Bentuk: Praktikum/PBL. Metode: hands-on practice, mentoring, code review, demo, dan presentasi. & 1 x 2 x 50' tatap muka; 1 x 2 x 50' tugas/praktik mandiri. & Mahasiswa mengerjakan aktivitas terarah untuk form dan validasi dan menunjukkan evidence hasil belajar. & Bentuk penilaian: Form dan validasi. Mengacu rubrik RTM. & Ketepatan konsep; kualitas artefak/praktik; validasi dan dokumentasi. & 5 \\
7 & SCPMK0603-03402 & Authentication dan authorization. & Modalitas: Blended Learning. Bentuk: Praktikum/PBL. Metode: hands-on practice, mentoring, code review, demo, dan presentasi. & 1 x 2 x 50' tatap muka; 1 x 2 x 50' tugas/praktik mandiri. & Mahasiswa mengerjakan aktivitas terarah untuk authentication dan authorization dan menunjukkan evidence hasil belajar. & Bentuk penilaian: Authentication dan authorization. Mengacu rubrik RTM. & Ketepatan konsep; kualitas artefak/praktik; validasi dan dokumentasi. & 5 \\
8 & SCPMK0209-03401, SCPMK0603-03402 & UTS modul CRUD Laravel. & Modalitas: Blended Learning. Bentuk: Praktikum/PBL. Metode: hands-on practice, mentoring, code review, demo, dan presentasi. & 1 x 2 x 50' tatap muka; 1 x 2 x 50' tugas/praktik mandiri. & Mahasiswa mengerjakan aktivitas terarah untuk uts modul crud laravel dan menunjukkan evidence hasil belajar. & Bentuk penilaian: UTS modul CRUD Laravel. Mengacu rubrik RTM. & Ketepatan konsep; kualitas artefak/praktik; validasi dan dokumentasi. & 20 \\
9 & SCPMK0603-03402 & File import dan export. & Modalitas: Blended Learning. Bentuk: Praktikum/PBL. Metode: hands-on practice, mentoring, code review, demo, dan presentasi. & 1 x 2 x 50' tatap muka; 1 x 2 x 50' tugas/praktik mandiri. & Mahasiswa mengerjakan aktivitas terarah untuk file import dan export dan menunjukkan evidence hasil belajar. & Bentuk penilaian: File import dan export. Mengacu rubrik RTM. & Ketepatan konsep; kualitas artefak/praktik; validasi dan dokumentasi. & 5 \\
10 & SCPMK0603-03402 & API. & Modalitas: Blended Learning. Bentuk: Praktikum/PBL. Metode: hands-on practice, mentoring, code review, demo, dan presentasi. & 1 x 2 x 50' tatap muka; 1 x 2 x 50' tugas/praktik mandiri. & Mahasiswa mengerjakan aktivitas terarah untuk api dan menunjukkan evidence hasil belajar. & Bentuk penilaian: API. Mengacu rubrik RTM. & Ketepatan konsep; kualitas artefak/praktik; validasi dan dokumentasi. & 5 \\
11 & SCPMK1006-03403 & Testing dan debugging. & Modalitas: Blended Learning. Bentuk: Praktikum/PBL. Metode: hands-on practice, mentoring, code review, demo, dan presentasi. & 1 x 2 x 50' tatap muka; 1 x 2 x 50' tugas/praktik mandiri. & Mahasiswa mengerjakan aktivitas terarah untuk testing dan debugging dan menunjukkan evidence hasil belajar. & Bentuk penilaian: Testing dan debugging. Mengacu rubrik RTM. & Ketepatan konsep; kualitas artefak/praktik; validasi dan dokumentasi. & 5 \\
12 & SCPMK1006-03403 & Kuis/review kualitas Laravel. & Modalitas: Blended Learning. Bentuk: Praktikum/PBL. Metode: hands-on practice, mentoring, code review, demo, dan presentasi. & 1 x 2 x 50' tatap muka; 1 x 2 x 50' tugas/praktik mandiri. & Mahasiswa mengerjakan aktivitas terarah untuk kuis/review kualitas laravel dan menunjukkan evidence hasil belajar. & Bentuk penilaian: Kuis/review kualitas Laravel. Mengacu rubrik RTM. & Ketepatan konsep; kualitas artefak/praktik; validasi dan dokumentasi. & 5 \\
13 & SCPMK0603-03402 & Integrasi proyek PBL. & Modalitas: Blended Learning. Bentuk: Praktikum/PBL. Metode: hands-on practice, mentoring, code review, demo, dan presentasi. & 1 x 2 x 50' tatap muka; 1 x 2 x 50' tugas/praktik mandiri. & Mahasiswa mengerjakan aktivitas terarah untuk integrasi proyek pbl dan menunjukkan evidence hasil belajar. & Bentuk penilaian: Integrasi proyek PBL. Mengacu rubrik RTM. & Ketepatan konsep; kualitas artefak/praktik; validasi dan dokumentasi. & 2 \\
14 & SCPMK1006-03403 & Review kode dan refactoring. & Modalitas: Blended Learning. Bentuk: Praktikum/PBL. Metode: hands-on practice, mentoring, code review, demo, dan presentasi. & 1 x 2 x 50' tatap muka; 1 x 2 x 50' tugas/praktik mandiri. & Mahasiswa mengerjakan aktivitas terarah untuk review kode dan refactoring dan menunjukkan evidence hasil belajar. & Bentuk penilaian: Review kode dan refactoring. Mengacu rubrik RTM. & Ketepatan konsep; kualitas artefak/praktik; validasi dan dokumentasi. & 2 \\
15 & SCPMK1006-03403 & Deployment preparation. & Modalitas: Blended Learning. Bentuk: Praktikum/PBL. Metode: hands-on practice, mentoring, code review, demo, dan presentasi. & 1 x 2 x 50' tatap muka; 1 x 2 x 50' tugas/praktik mandiri. & Mahasiswa mengerjakan aktivitas terarah untuk deployment preparation dan menunjukkan evidence hasil belajar. & Bentuk penilaian: Deployment preparation. Mengacu rubrik RTM. & Ketepatan konsep; kualitas artefak/praktik; validasi dan dokumentasi. & 1 \\
16 & SCPMK0209-03401--SCPMK1006-03403 & PBL proyek Laravel. & Modalitas: Blended Learning. Bentuk: Praktikum/PBL. Metode: hands-on practice, mentoring, code review, demo, dan presentasi. & 1 x 2 x 50' tatap muka; 1 x 2 x 50' tugas/praktik mandiri. & Mahasiswa mengerjakan aktivitas terarah untuk pbl proyek laravel dan menunjukkan evidence hasil belajar. & Bentuk penilaian: PBL proyek Laravel. Mengacu rubrik RTM. & Ketepatan konsep; kualitas artefak/praktik; validasi dan dokumentasi. & 10 \\
17 & SCPMK1006-03403 & UAS validasi proyek Laravel. & Modalitas: Blended Learning. Bentuk: Praktikum/PBL. Metode: hands-on practice, mentoring, code review, demo, dan presentasi. & 1 x 2 x 50' tatap muka; 1 x 2 x 50' tugas/praktik mandiri. & Mahasiswa mengerjakan aktivitas terarah untuk uas validasi proyek laravel dan menunjukkan evidence hasil belajar. & Bentuk penilaian: UAS validasi proyek Laravel. Mengacu rubrik RTM. & Ketepatan konsep; kualitas artefak/praktik; validasi dan dokumentasi. & 10 \\
\end{longtblr}
